\section{Discussion}
\label{sec:discussion}

\paragraph{What the guarantee says.}
With probability $1-\delta$ over the calibration sample, the population risk of the deployed policy, accepted and
substituted answers together, is at most $\alpha$. It says nothing about a single item, about a finite test fold, about
a distribution other than the one the units were drawn from, or about the coverage that results. A sampling model is
part of it: units are i.i.d., and when a benchmark is split at random the iid assumption is a modelling choice about the population
the benchmark represents.

\paragraph{What the evidence supports.}
Under a protocol in which the proved procedure is the evaluated procedure, certified repair gains between $0.8$ and $7.9$ points
of coverage over certified filtering when the gate is fixed on external data; the sign is supported by $65$--$99\%$ of
image-bootstrap datasets and the interval excludes zero on one cell. The statement a reader can rely on is therefore
that the gain is positive in expectation on these benchmarks and small, and that more images would be needed to
size it. Two of the three mechanisms the earlier analysis offered survive in weaker form. Repair does add a few points of coverage; it
does so mostly by confirming answers the model already gave, and corrects almost none of the model's false positives.
The correctness score dominates the coverage, and the start of the fixed sequence dominates the repair on small folds.

\paragraph{Limitations.}
(i) The task is binary object existence, with one detector, three backbones and two benchmark families whose
images overlap; $642$ distinct images are involved in total. (ii) The channel-2 operating point $\theta=0.15$ was fixed
in the extraction scripts and is not varied here; we do not claim it was chosen independently of the benchmarks.
(iii) The external gate for the POPE cells was chosen on AMBER, whose image sources we could not verify to be disjoint from
POPE's COCO images; the guarantee requires the gate to be independent of the calibration sample, and any image overlap
would weaken that independence slightly. (iv) The i.i.d.\ assumption is violated within images when errors are correlated; here the
measured correlation is slightly negative, and \S\ref{sec:validity-audit} shows the guarantee degrades with positive
correlation unless calibration uses one item per image. (v) The bootstrap interval reflects sampling of images from the benchmark, not
of benchmarks from the world. (vi) Free-form generation is not covered and mid-sequence repair worsens the continuation.
(vii) We did not rerun the VLMs, extract signals on additional images, or run BCEA or CEBC.

\paragraph{Next steps.}
The measurement that would most improve the paper is more images: processing the remaining POPE images and a
second detector would narrow the intervals the bootstrap now shows to be wide. A repair rule that
respects the generation context, tested on the free-form setting with the matched-prefix design, is the natural
extension.

\section{Conclusion}
\label{sec:conclusion}

Certified correction replaces an answer the model would have abstained on with a detector's answer and
certifies the emitted set as a whole; the procedure is valid under i.i.d.\ sampling when thresholds and gate are
fixed before calibration. Under a protocol built on that requirement the coverage gain over filtering is small and positive, with an
uncertainty that the available images do not resolve. Most repaired items re-emit the model's answer with detector
confirmation, the model's false positives are not corrected, and the coverage comes mainly from the correctness score.
In generation, mid-sequence repair worsens the continuation. The reproducible pipeline and all negative results are released.
